% !TEX program = xelatex
% Compile: xelatex -interaction=nonstopmode -halt-on-error open-cake-ir-technical-report.tex
% This single source intentionally has no external figure or bibliography files.
\documentclass[UTF8,11pt,a4paper,fontset=none]{ctexart}
\setCJKmainfont{Songti SC}[BoldFont=Hiragino Sans GB]
\setCJKsansfont{Hiragino Sans GB}[BoldFont=Hiragino Sans GB]
\setCJKmonofont{Songti SC}
\setCJKfamilyfont{zhhei}{Hiragino Sans GB}
\usepackage[margin=25mm,headheight=15pt]{geometry}
\usepackage{amsmath,amssymb,booktabs,tabularx,array,enumitem,xcolor}
\usepackage{hyperref,fancyhdr}
\definecolor{ink}{HTML}{17324D}
\definecolor{muted}{HTML}{526579}
\definecolor{accent}{HTML}{176B80}
\hypersetup{colorlinks=true,linkcolor=ink,urlcolor=accent,citecolor=accent,pdftitle={open-cake-ir: Agent 驱动的 GPU Kernel 与编译器协同演进技术报告},pdfauthor={Haiyan Qin}}
\setlength{\parindent}{2em}
\setlength{\parskip}{0.28em}
\renewcommand{\arraystretch}{1.25}
\setlist[itemize]{leftmargin=2.3em,itemsep=0.15em,topsep=0.25em}
\setlist[enumerate]{leftmargin=2.3em,itemsep=0.15em,topsep=0.25em}
\pagestyle{fancy}
\fancyhf{}
\fancyhead[L]{\small\color{muted} open-cake-ir 技术报告}
\fancyhead[R]{\small\color{muted} 2026-09-23}
\fancyfoot[C]{\thepage}
\newcommand{\code}[1]{\texttt{#1}}
\newcommand{\repo}[2]{\href{https://github.com/qhy991/open-cake-ir/blob/0fe3a4476765a1d323e9d114d1c9af2f324dd767/#1}{#2}}
\newcommand{\scope}[1]{\textcolor{accent}{\textbf{#1}}}
\newcolumntype{Y}{>{\raggedright\arraybackslash}X}

\begin{document}
\begin{titlepage}
\centering
\vspace*{23mm}
{\Large\color{accent} open-cake-ir\par}
\vspace{8mm}
{\Huge\bfseries Agent 驱动的 GPU Kernel\\[3mm]与编译器协同演进\par}
\vspace{5mm}
{\Large 技术报告\par}
\vspace{9mm}
{\large\color{muted} A Technical Report on Agent-Driven GPU Kernel\\ and Compiler Co-Evolution\par}
\vspace{17mm}
{\Large 秦海岩\quad Haiyan Qin\par}
\vspace{3mm}
{\normalsize 北京航空航天大学\par}
\vfill
\begin{minipage}{0.86\textwidth}
\small
\textbf{版本与证据边界。}本报告整理仓库在 2026 年 9 月 23 日的实现与已发布证据；正文所读源码固定为
\code{0fe3a4476765a1d323e9d114d1c9af2f324dd767}。历史设备实验仍绑定各自的源码、精确目标、
Workload、oracle 和计时协议，不能因本报告的新版本而重新归属。文内给出固定提交的来源链接。
\end{minipage}\par
\vspace{8mm}
{\large 2026 年 9 月 23 日\par}
\end{titlepage}

\begin{abstract}
GPU Kernel Agent 可以快速提出候选实现，但候选为何非法、何处消耗资源、哪次设备测量可以支持收益，
往往分散在代码生成器、运行时和实验记录中。open-cake-ir 将单 Kernel 的硬件执行决策表示为带类型的
\emph{Schedule}，以 \emph{Program} 组合阶段；Compiler 在生成目标源码前给出局部化诊断，Lab 在冻结的
Workload 和 Run 合同下搜索候选，Evaluation 以外部 oracle、目标计时和 profiler 检查结果，Evidence
保存可审计的观察。编译器缺口由运行之外的后继提交修复，并通过语料门验证，形成 Kernel 优化与
编译能力演进的两个循环。

本文给出系统的职责、表示、检查和实验协议，并将证据按层级陈述。当前源码的 179 项 Corpus Gate
用例符合预期；B300、C550 等目标已有范围明确的设备案例，包括合格的局部改进与未通过计时质量门的
完整 Program。现有记录支持平台接入和若干固定任务的局部优化；跨硬件经验材料或显式变换能否降低
Agent 搜索成本，仍需按预注册的四组处理实验验证。本文不将软件合同、设备正确性和因果收益混为一项结论。
\end{abstract}

\noindent\textbf{关键词：}GPU Kernel；中间表示；静态验证；Agent；实验可复查性；跨硬件优化

\vspace{1em}
\noindent\textbf{English abstract.}
open-cake-ir is a research system in which agents author typed, hardware-explicit kernel schedules,
receive localized compiler feedback, and validate candidates against external oracles and target-specific
measurements. The compiler evolves between frozen optimization runs, with corpus-gated changes. This
report separates software contract coverage, device correctness, qualified local performance results,
and the still untested causal claim of cross-hardware knowledge transfer. Its implementation snapshot is
commit \code{0fe3a447}; each historical measurement retains its own source and assay boundary.

\clearpage
\tableofcontents
\clearpage

\section{问题、研究对象与论点}

给定同一数学算子，GPU 实现仍须决定任务分块、线程分工、存储层次、访问坐标、同步与执行顺序。
一个只输出延迟的搜索过程很难说明某次失败属于候选调度、编译器表达能力、硬件准入，还是测量噪声。
CAKE 论文提出让 Agent 编辑显式执行计划，并以编译器分析和外部评测形成反馈闭环\cite{cake}。
open-cake-ir 是对这一研究路线的独立工程实现；论文报告的模型实验和速度数字不作为本仓库结果\cite{papercontract}。

本报告的研究对象是\emph{可复查的 Agent 编写环境}：一份候选包含什么决策，系统能在什么范围内拒绝，
设备证据如何决定接受，以及真实缺口怎样改变下一版 Compiler。核心问题可以表述为：
\begin{quote}
在保持 Workload 语义和实验条件可追溯的前提下，怎样让 Agent 的 Kernel 决策可检查、
让失败有负责位置，并让重复经验进入后继编译能力？
\end{quote}
本文主要陈述三项已实现的机制：\textbf{显式的 Program/Schedule 合同}、\textbf{分层诊断与目标精确准入}、
\textbf{冻结运行和追加式证据}。跨硬件知识迁移的增量价值是后续实验问题，不是现有软件实现的推论。

\section{系统合同与职责边界}

\subsection{从任务到可审计结论}
Workload Contract 固定数学语义、输入域、外部 oracle 和容差。Schedule 描述一个 Kernel 的资源、
执行组、操作、地址与同步；Program 再声明公共张量、各阶段顺序和绑定。RunSpecification 固定一次
优化所用的源码、目标、可见材料、工具权限和预算。StudyPlan 仅在研究对照时进一步固定处理、分配、
终点与分析。这些对象的职责不同，不能用一个候选的局部通过来代替 Study 结论\cite{architecture,study}。

\begin{table}[htbp]
\centering\small
\caption{一次候选从表达走向结论时的负责位置。箭头表示依赖方向，不表示某层的通过自动推出下一层的通过。}
\begin{tabularx}{\textwidth}{@{}p{25mm}Y Y@{}}
\toprule
对象或模块 & 拥有的事实 & 对外给出的产物 \\
\midrule
Workload & 数学、形状、输入分布、外部答案 & 不随候选改变的验收合同 \\
Compiler & IR 类型、Target、合法性、后端准入与生成 & Assessment、Finding、目标源码 \\
Lab & 材料、候选搜索、预算与 Run 控制 & 已提交候选和过程事件 \\
Evaluation & 外部正确性、目标计时质量和剖析 & 独立的检查、计时和 profiler 收据 \\
Evidence & 原始文件、事件、回放与审计 & 可由原始材料重建的报告 \\
Study & 处理分配、估计量和缺失策略 & 限定研究问题的分析结果 \\
\bottomrule
\end{tabularx}
\end{table}

Compiler 是独立产品核心，不导入 Lab、provider、Workload 实现或证据库。Lab 调用 Compiler 与
Evaluation，Evidence 保存其观察。这样的依赖关系使同一个 Compiler 可用于离线检查，也防止
实验策略偷偷进入 Schedule 的语义判断。系统源码与边界见\repo{docs/ARCHITECTURE.md}{架构章节}和
\repo{CONTEXT-MAP.md}{职责地图}。

\subsection{两个时钟：候选搜索与编译器演进}
内循环在\emph{冻结的} Compiler、Workload、Target、oracle 和测量合同下进行：生成结构不同的候选，
先通过构造检查、Verifier 和后端预检筛选，再对幸存者做设备正确性、计时与剖析，最后把诊断
送回下一次候选。外循环由累积的 Finding 触发：若问题属于通用合法性、缺失原语、lowering 或
目标校准，维护者在后继提交中修改其负责模块，验证后才开启新 Campaign。一次运行不能边测边
改变尺子\cite{architecture,acceptance}。

\section{表示：Workload、Program 与 Schedule}

\subsection{显式的计划，而非新的布局代数}
受限 Python 前端与 JSON 文档形式进入同一规范的 typed Schedule；Python 入口解析受支持的语法，
并不执行任意宿主 Python 程序。Schedule 记录执行组、缓冲区、访问映射、循环、操作与同步，
性能相关的硬件选择因此可被检查。它没有把 layout 建成一阶代数；对目标硬件的具体存储与访问
承诺由相应 Target 和后端解释并验证。Triton 路线仍把部分物理 placement、寄存器分配与指令
选择交给下游工具链，不能把其生成结果当作 Schedule 已经精确指定的事实\cite{irguide,architecture}。

例如对一行 Softmax，Workload 给出 $y_j=\exp(x_j-m)/\sum_k\exp(x_k-m)$，其中
$m=\max_k x_k$；候选 Schedule 另选每个执行组处理的行、列块、归约方式与写回坐标。
若读取、归约和写回对同一行的坐标约定不一致，问题属于数据一致性；若两个执行组写同一输出，
问题属于程序安全；若目标不提供所需指令，则属于硬件符合性或后端资格。数学式相同不能消除
这些执行计划差异。

Program 负责多个 Schedule 的公共输入输出、阶段顺序与张量绑定。完整 Program 的正确性要检查
所有阶段的最终输出；单 Kernel 的计时或剖析不能自然推广为 Program、模型或服务收益。

\subsection{Target 是精确的硬件文档}
当前源码声明八个精确目标，分属 NVIDIA、Apple、AMD、Hygon 和 MetaX 五个 Vendor。
Target 文档拥有自己的 lane 宽度、存储空间、指令合同、代码对象及资源事实；未声明事实给出
“未建模”的报告，不借用相似目标的常数。Backend 是生成机制而不是 Vendor：现有
\code{triton}、\code{metal}、\code{native\_cuda} 和 \code{cutlass\_cute\_dsl} 各声明可生成的
代码对象。Compiler 对目标和后端做精确匹配，不能把未声明目标降级到邻近架构\cite{targets,architecture}。

\subsection{表示变更的验收义务}
新增 IR 原语须同时给出类型、效应、合法性、静态分析和目标行为；只有语法而没有这些语义，
不能成为已准入能力。仓库将易编写、性能决策可见、规范形式、构造时类型检查、可分析、
语料门、分析一致与硬件根据作为表示的八项设计约束\cite{cake,irguide}。某算子若能由现有原语
组合实现，应先添加 Workload Contract；只有真实表达缺口才要求扩大 Compiler。

\section{Compiler：局部化检查到目标源码}

\subsection{分层拒绝与反馈}
Compiler 首先解析与构造类型正确的 IR，然后以四类合同检查：\code{schedule\_semantics}、
\code{hardware\_conformance}、\code{data\_consistency} 和 \code{program\_safety}。
Assessment 中的 Finding 指向字段或源码区域，带稳定代码、类别、严重度和阻断范围；
非阻断的 guidance 与拒绝原因分开。后端预检负责该生成路线特有的能力边界。静态分析只能
在已建模的范围内阻断；未声明的资源上限和下游产生的物理资源需求需要另行观察\cite{architecture,acceptance}。

一个可复查的小反例是 FMA 候选漏掉第三个输入：构造仍指向同一数学操作，Assessment 在
\code{operations[3].reads} 报 \code{ELEMENTWISE\_ARITY}，指出该操作需要三个输入。
作者因此可以修复具体数据边，无需先耗费设备时间。通过静态检查只允许候选继续；
它不证明工具链编译、GPU 数值正确或速度收益\cite{architecture}。

\subsection{源码生成与下游工具链}
合法 Schedule 经选定后端输出可检查的目标源码。Triton 路线再通过厂商对应的 Triton 工具链
编译；Metal、原生 CUDA 与 CuTe DSL 路线各保留自己的预检与生成条件。Target 声明的代码对象
随生成路线进入工具链和 Executor，不能由目标名字猜测架构。生成源码、编译产物、设备可加载性、
外部正确性与实测延迟是五个递进但互不等同的检查点\cite{architecture}。

\subsection{Corpus Gate 的含义}
固定源码 \code{0fe3a447} 的状态页报告 \textbf{179/179} 项语料检查符合已审查预期；
其中 Apple family7 与 family9 是已声明但未被该语料检查的目标。
该数值证明当前软件合同与预期相符，不证明这些目标的设备运行或性能。语料预期必须作为独立
审查决定维护，不能为使 Gate 变绿而在验证前重生成\cite{status,acceptance}。

\section{实验方法与证据等级}

\subsection{一次 Run 的四阶段}
系统按下列顺序运行候选搜索：
\begin{enumerate}
\item 生成结构不同的 Schedule 或 Program 候选，并记录具体变更；
\item 以 IR 构造、Verifier、backend preflight 和有目标覆盖的成本估计筛选，未覆盖时明确弃权；
\item 对幸存者使用外部 oracle 检查完整输出，再在目标计时器、状态重置与质量门下测量，并保留 profiler 观察；
\item 将失败归于候选、Verifier、成本模型、IR 或 lowering，决定继续搜索或登记后继 Finding。
\end{enumerate}
成本估计用于设备时间之前的排序与过滤，不能决定最终接受。每个目标分别声明计时器、测量区间
和样本前设备状态；缺少其中一项时应报告测量覆盖限制，而不是继承其他目标的延迟语义。

\subsection{研究对照的单位和参考访问}
正式比较的实验单位是独立 Run，不是 Turn、候选或检查点。干净起点只能看到规格、oracle、
高层代码和被授权的参考资料，不能读取目标的完整低层实现；已知 Kernel 复现则可以读取该实现。
直接低层编写与 Cake 编写是 Authoring Environment 的不同处理条件，需要共同的目标、任务、
模型、预算与评测协议。搜索内的最佳读数须经候选固定后的独立确认；基础设施故障应保留为缺失，
不能重写成数值错误\cite{papercontract,study}。

\subsection{结论阶梯}
表\ref{tab:evidence-ladder} 是本文解读每个案例的最小证据门。任一层的通过只授予该层的陈述资格。
\begin{table}[htbp]
\centering\small
\caption{结论阶梯；“通过”不沿表格自动向下传递。}
\label{tab:evidence-ladder}
\begin{tabularx}{\textwidth}{@{}p{31mm}Y Y@{}}
\toprule
可陈述的结果 & 必要证据 & 尚需另行验收 \\
\midrule
结构与目标准入 & typed IR、精确 Target、Finding 与后端预检 & 工具链编译和加载 \\
设备正确性 & 完整输出对外部 oracle，覆盖规定形状与分布 & 计时质量和性能收益 \\
局部性能 & 固定候选与基线、共同计时区间、质量通过 & 跨任务或跨硬件泛化 \\
Program 或服务收益 & 完整调用路径、dispatch 与应用层验收 & Agent 处理的因果效果 \\
研究归因 & 预注册处理、独立 Run、共同确认与缺失规则 & 未测试任务与平台的外推 \\
\bottomrule
\end{tabularx}
\end{table}

Evidence 保存候选、原始样本和事件；报告是由这些记录派生的视图。源码身份由干净 Git 提交
决定，Compiler 为 \code{open-cake-ir@<commit>}，Executor 为 \code{<target>@<commit>}。
本报告的阅读提交只界定当前叙述所见的软件，不能代替历史实验的执行提交或文件监管条件。

\section{实现覆盖与代表性设备案例}

\subsection{平台覆盖是分层事实}
表\ref{tab:platforms} 概括仓库已发布的代表性证据。它不是跨厂商性能排行；同一平台的不同
Workload 也可能有不同的计时和确认资格。完整细节由各平台专题与原始收据负责\cite{architecture}。
\begin{table}[htbp]
\centering\small
\caption{截至报告快照的代表性平台能力与主要解释边界。}
\label{tab:platforms}
\begin{tabularx}{\textwidth}{@{}p{23mm}Y Y@{}}
\toprule
平台 & 已报告的执行范围 & 需保留的边界 \\
\midrule
NVIDIA & B300 的单 Kernel、部分完整 Program；B200 另有固定案例 & FlashInfer 外部比较有通过与未通过质量门的边，按任务解读 \\
Apple & Metal 固定算子与 TaskLab 优化记录 & 已验收设备和计时口径不推广到全部 family 或完整 Program \\
AMD & \code{gfx1151} 设备调查与 smoke & 计时器绝对时间未对齐，尚无合格的完整 Agent 收益结论 \\
Hygon & BW1101 / \code{gfx938} 固定任务与 HIP 路线 & 短 Kernel 的计时分辨能力限制部分比较 \\
MetaX & C550 固定单 Kernel 及若干完整输出；单 Kernel 配对计时和 profiler & 完整 Program 普通 Run 计时与完整 provider 闭环仍需验收 \\
\bottomrule
\end{tabularx}
\end{table}

\subsection{案例 A：C550 的局部 tiling 收益}
在 C550 的 FP16 GEMM 固定形状 $M=17,N=128,K=2048$ 上，已验收基线保持不变，
显式地把 M tile 从 64 改为 32。独立确认的基线和候选中位数分别为
\textbf{72.448 与 58.368\,$\mu$s}，比值 \textbf{1.241 倍}；正确性 case 与测量质量门通过。
这是 \code{local\_serialized} 范围内的本机 authoring 对照，支持这个形状上的局部调度改进。
它没有比较“向 Agent 提供 NVIDIA 经验”与“不提供经验”的处理差异，也不是完整 Program 的
自动优化 Run\cite{metax}。同一专题中的 17 个 Workload 变体、86/86 个完整输出 case
属于正确性证据；其计时为缺失，不能拿来扩大上述收益范围。

\subsection{案例 B：B300 的对齐条件与质量门}
B300 FlashInfer 任务 026 的守卫式对齐候选，相对同次冻结的 generic \code{sliced-w8}
对照得到合格的 \textbf{1.081 倍}局部收益；比较所用中位数为 6.4005 与 5.920\,$\mu$s。
同一候选对外部参考的边没有通过 CV 质量门，因此不能从该案例推出对外部参考的合格胜负。
另一任务 008 的完整两阶段 Program 已通过 2550 份数值快照，但新候选相对对照及外部参考
的计时边均未通过质量门：它证明完整程序可运行和算对，不提供合格的性能结论\cite{nvidia}。

\subsection{案例的合取边界}
两个局部改进来自不同硬件、Workload、源码与计时机制，不能平均，也不能作为一个跨硬件
机制已迁移的轨迹。负结果、\code{close\_null}、未完成和质量失败与正结果同样保留。
截至此报告所读的\repo{reports/current/STATUS.md}{当前状态页}，项目没有统一的已验收 Study Report
索引；本文不计算全项目成功率或总加速比。

\section{跨硬件知识迁移：机制与待验证估计量}

\subsection{可复用的单位与负责位置}
候选优化可能包含两种不同的知识：一个带适用条件的\emph{语义改写}，以及在某个硬件上选出的
tile、执行组或 pipeline 参数。前者可在 IR 中形成显式 pass，声明前提并返回完整候选或拒绝原因；
后者要在目标 Target 上重新选择和验证。以私有中间张量的融合为例，数据流私有性、舍入、消费者
和生命周期是机制条件；寄存器占用与具体收益属于目标设备事实。来源设备的速度比不随机制迁移
\cite{transfer}。当一次研究没有可推广改写时，记录 \code{No promotion} 是正常结果。

\subsection{E/P 四组实验}
仓库已实现对经验材料 $E$ 与显式变换调用权 $P$ 的四格隔离和软件协议，但尚没有共同目标、
基线与预算下的实机迁移收益实验。四组分别为 $E0P0$、$E1P0$、$E0P1$ 和 $E1P1$；
$P1$ 的最小接口合同本身会暴露必要机制信息，因而 $E$ 的效果只指额外说明与案例。
各组共享 typed IR、Target、oracle、计时质量门和可表达的基础能力。材料与调用权限须通过
真实隔离执行，不能只由提示词声称关闭\cite{study}。

令 $p_{ep}$ 为在预注册预算内经新鲜确认、达到本机预定改善阈值的 Run 比例。预定主要对比为
\begin{align*}
\Delta_E&=\tfrac12[(p_{10}-p_{00})+(p_{11}-p_{01})],\\
\Delta_P&=\tfrac12[(p_{01}-p_{00})+(p_{11}-p_{10})],\\
\Delta_{EP}&=p_{11}-p_{10}-p_{01}+p_{00}.
\end{align*}
独立 Run 是重复单位；正式多任务研究以任务等权汇总并保留组内配对结构。搜索成本、失败、
预算耗尽和基础设施缺失须按预注册规则保留。协议的存在证明问题可被提出和执行，只有冻结
处理后的设备实验才能估计上述三个量。

\section{限制、可复查性与结论}

本文的最强软件结论是当前源码上的表示、检查、生成和控制协议，以及 179 项语料预期相符。
最强设备结论是指定源码、目标与固定 Workload 上的正确性和局部合格改善。当前材料不足以
估计跨硬件经验迁移的因果收益、普遍性能优势，或模型与服务层面的端到端收益。
静态分析的覆盖也受目标声明和后端实际资源分配限制；缺乏目标计时校准时应报告未知，
而不是给出没有根据的延迟预测。

后续工作以同一职责边界推进：补齐各平台完整 Program 与目标框架验收；在有原始证据的
目标上校准测量与成本估计；将重复的失败形成有反例的 Verifier 规则、IR 能力或 guarded pass；
最后用预注册 E/P 实验检验迁移增量。这个顺序让 Kernel 结果、Compiler 演进与研究归因各有
自己的证据和版本。

\section*{阅读与引用说明}
\addcontentsline{toc}{section}{阅读与引用说明}
本文是作者维护的仓库技术报告，无 DOI、期刊或同行评审声明。引用设计时请固定本报告及其
所读源码提交；引用性能时还需追溯案例自己的执行源码、精确目标、Workload、baseline、
oracle、计时区间与质量门。以下文献中的仓库链接均固定到本文所读的 \code{0fe3a447}，
而实验原始收据的身份由各专题继续给出。

\begin{thebibliography}{99}
\bibitem{cake} Z. Ye et al. \emph{CAKE: Compiler--Agent Co-Design for Frontier Kernel Evolution}.
arXiv:2608.12629v1, 2026. \url{https://arxiv.org/html/2608.12629v1}.
\bibitem{architecture} H. Qin. \repo{docs/ARCHITECTURE.md}{系统架构与平台能力}, open-cake-ir, fixed source commit.
\bibitem{irguide} H. Qin. \repo{docs/IR_GUIDE.md}{IR 指南}, open-cake-ir, fixed source commit.
\bibitem{targets} H. Qin. \repo{compiler/revision.json}{Compiler 清单}与
\repo{compiler/targets/sm_103a.json}{Target 文档示例}, open-cake-ir, fixed source commit.
\bibitem{acceptance} H. Qin. \repo{docs/ACCEPTANCE_GATES.md}{验收门}, open-cake-ir, fixed source commit.
\bibitem{papercontract} H. Qin. \repo{docs/PAPER_CONTRACT.md}{论文与本地证据合同}, open-cake-ir, fixed source commit.
\bibitem{status} H. Qin. \repo{reports/current/STATUS.md}{当前状态}, open-cake-ir, fixed source commit.
\bibitem{nvidia} H. Qin. \repo{docs/results/nvidia/FLASHINFER_STATUS.md}{B300 FlashInfer 任务综述}, open-cake-ir, fixed source commit.
\bibitem{metax} H. Qin. \repo{docs/metax-c550.md}{C550 专题与 M17 独立确认}, open-cake-ir, fixed source commit.
\bibitem{transfer} H. Qin. \repo{docs/OPTIMIZATION_TRANSFER.md}{跨硬件优化知识迁移机制}, open-cake-ir, fixed source commit.
\bibitem{study} H. Qin. \repo{docs/OPTIMIZATION_TRANSFER_ABLATION.md}{E/P 消融实验设计}, open-cake-ir, fixed source commit.
\end{thebibliography}

\end{document}
